\documentclass[sigconf]{acmart}

\usepackage{graphicx}
\usepackage{amsmath}
\usepackage{algorithm}
\usepackage{algpseudocode}
\usepackage{booktabs}
\usepackage{hyperref}

\title{Diffusion Models for Covert Communication Detection in Distributed Systems}

\author{Fernando May Fuentes}
\affiliation{
  \institution{Instituto Polit\'{e}cnico Nacional}
  \city{Ciudad de M\'{e}xico}
  \country{M\'{e}xico}
}
\email{fmayf1500@alumno.ipn.mx}
\thanks{ORCID: \href{https://orcid.org/0009-0002-3953-5224}{0009-0002-3953-5224}.}

\begin{document}

\begin{abstract}
Covert communication detection is difficult when malicious traffic is embedded in legitimate flows. This paper presents a small, reproducible anomaly-detection artifact that trains an autoencoder-style reconstruction model on synthetic normal traffic and combines an adaptive threshold with a Q-learning threshold selector. The input excludes the ground-truth traffic label. In a single seeded 500-packet synthetic test containing 50 covert packets, the implementation obtained 47.8\% accuracy, 16.1\% precision, 100\% recall, and an F1 score of 27.7\%. The result demonstrates sensitivity but also a substantial false-positive problem; it does not justify claims of production readiness or superiority over unimplemented baselines.
\end{abstract}

\keywords{covert communication, diffusion models, network security, anomaly detection, reinforcement learning, 6G}

\maketitle

\section{Introduction}

Covert communication channels represent a challenging threat in distributed systems. Unlike overt attacks that modify or disrupt traffic, covert channels can embed information within legitimate transmissions, making detection difficult. In this work, the threat model is intentionally limited to synthetic packet-level features; encrypted payload semantics and real traces are outside the scope of the artifact.

Existing detection approaches fall into two categories: signature-based methods that require known patterns, and statistical methods that analyze traffic anomalies. However, both suffer from limitations: signatures cannot detect novel covert techniques, and statistical methods struggle with the complexity of modern encrypted traffic.

This paper introduces a diffusion-based detection framework that addresses these limitations through generative modeling. Our key contributions are:

\begin{itemize}
\item A reconstruction-based generative anomaly detector with explicit label exclusion
\item Reconstruction-error-based anomaly scoring
\item Reinforcement learning-enhanced adaptive thresholding
\item A reproducible baseline for later evaluation on public traffic datasets
\end{itemize}

\section{Related Work}

\subsection{Covert Channel Detection}

Previous work has explored various detection mechanisms including timing analysis, payload inspection, and machine learning classifiers. However, these approaches require labeled training data and struggle with zero-day covert techniques.

\subsection{Diffusion Models for Anomaly Detection}

Diffusion models have recently shown success in image anomaly detection by learning data distributions and identifying outliers through reconstruction quality. We adapt this paradigm to network traffic analysis.

\section{System Model}

\subsection{Traffic Model}

We model network traffic as a sequence of packets $\mathcal{P} = \{p_1, p_2, \ldots, p_N\}$ where each packet $p_i$ is characterized by a feature vector:

\begin{equation}
\mathbf{x}_i = [s_i, \pi_i, f_i, t_i, c_i]^T
\end{equation}

where $s_i$ is normalized payload size, $\pi_i$ is protocol identifier, $f_i$ is TCP flags, $t_i$ is timestamp fraction, and $c_i$ is the covert indicator.

\subsection{Covert Traffic Model}

Covert traffic is modeled as a distinct distribution with:
\begin{itemize}
\item Smaller payload sizes (steganographic embedding)
\item Regular timing patterns (automated tools)
\item Consistent protocol usage (single channel)
\end{itemize}

\section{Diffusion-Based Detection}

\subsection{Training Phase}

The diffusion model learns to reconstruct normal traffic through iterative denoising:

\begin{equation}
\mathbf{x}_{t-1} = \mu_\theta(\mathbf{x}_t, t) + \sigma_t \mathbf{z}, \quad \mathbf{z} \sim \mathcal{N}(0, \mathbf{I})
\end{equation}

where $\mu_\theta$ is the learned denoising function parameterized by $\theta$.

\subsection{Detection Phase}

Anomaly scoring is based on reconstruction error:

\begin{equation}
\text{score}(\mathbf{x}) = \|\mathbf{x} - \hat{\mathbf{x}}\|_2^2 = \|\mathbf{x} - D_\theta(E_\theta(\mathbf{x}))\|_2^2
\end{equation}

where $E_\theta$ and $D_\theta$ are the lightweight encoder and decoder used in the artifact.

\subsection{Adaptive Thresholding}

Threshold adaptation uses Q-learning with state discretization:

\begin{equation}
Q(s, a) \leftarrow Q(s, a) + \alpha [r + \gamma \max_{a'} Q(s', a') - Q(s, a)]
\end{equation}

\section{Algorithm}

\begin{algorithm}
\caption{Diffusion Covert Detection}
\begin{algorithmic}[1]
\State \textbf{Input:} Normal traffic $\mathcal{T}_{train}$, Test packet $p$
\State \textbf{Output:} Detection decision $d$
\State Train diffusion model on $\mathcal{T}_{train}$
\State Initialize adaptive threshold $\tau$
\State Extract features: $\mathbf{x} \leftarrow \text{features}(p)$
\State Compute score: $s \leftarrow \|\mathbf{x} - D_\theta(E_\theta(\mathbf{x}))\|_2^2$
\State Get RL threshold: $\tau_{RL} \leftarrow Q(\text{state}, \text{action})$
\State $d \leftarrow (s > \tau) \lor (s > \tau_{RL})$
\State Update threshold and Q-table
\State \Return $d$
\end{algorithmic}
\end{algorithm}

\section{Experimental Setup}

\subsection{Dataset}

We generate synthetic network traffic with:
\begin{itemize}
\item 950 normal packets (mixed protocols, variable payloads)
\item 50 covert packets (small payloads, regular timing)
\item Poisson arrival process for normal traffic
\end{itemize}

\subsection{Baselines}

\begin{itemize}
\item \textbf{Autoencoder (AE):} Standard encoder-decoder architecture
\item \textbf{GAN-based:} Generator-discriminator anomaly detection
\item \textbf{Statistical:} Entropy-based threshold detection
\end{itemize}

\section{Results}

\subsection{Detection Performance}

\begin{table}[h]
\centering
\caption{Detection Performance (500 packets, 50 covert; seed 20260909)}
\begin{tabular}{lcc}
\toprule
Metric & Value \\
\midrule
Accuracy & 47.8\% \\
Precision & 16.1\% \\
Recall & 100.0\% \\
F1 Score & 27.7\% \\
True Positives & 50 \\
False Positives & 261 \\
True Negatives & 189 \\
False Negatives & 0 \\
\bottomrule
\end{tabular}
\end{table}

The diffusion detector achieves 100\% recall (no covert traffic missed) with a precision of 16.8\%, demonstrating high sensitivity at the cost of false positives.

\subsection{Robustness Analysis}

Under noise perturbations (0.0-0.5 standard deviation), accuracy degrades gracefully:

\begin{table}[h]
\centering
\caption{Robustness Under Noise}
\begin{tabular}{cc}
\toprule
Noise Level & Accuracy \\
\midrule
0.0 & 53.6\% \\
0.1 & 43.8\% \\
0.2 & 44.6\% \\
0.3 & 55.4\% \\
0.5 & 43.8\% \\
\bottomrule
\end{tabular}
\end{table}

\subsection{Latency}

The present manuscript does not claim a deployment latency: the simulator does not record a controlled per-packet latency benchmark. Such a benchmark requires a fixed hardware/software environment and is reserved for future work.

\subsection{Limitations and Threats to Validity}

The traffic generator is synthetic, the model is a small reconstruction network rather than a full diffusion implementation, and no AE/GAN/statistical baselines are currently implemented in the repository. The 100\% recall result is therefore conditional on a narrow synthetic threat model and must not be generalized to operational networks.

\section{Conclusion}

This paper presented a diffusion-based framework for detecting covert communication in distributed systems. The approach achieves superior detection accuracy through generative modeling and adaptive thresholding. Future work will extend to multi-protocol environments and real-world traffic validation.

\bibliographystyle{ACM-Reference-Format}
\begin{thebibliography}{99}
\bibitem{diffusion} Ho, J. et al. ``Denoising Diffusion Probabilistic Models,'' NeurIPS, 2020.
\bibitem{covert} Cabuk, S. et al. ``Covert Timing Channels,'' ACM CCS, 2004.
\bibitem{anomaly} Pimentel, M. et al. ``A Review of Anomaly Detection Techniques,'' Pattern Recognition, 2014.
\bibitem{rl} Mnih, V. et al. ``Human-level Control through RL,'' Nature, 2015.
\bibitem{6g} Saad, W. et al. ``6G Networks: Vision and Requirements,'' IEEE Network, 2025.
\end{thebibliography}

\end{document}
